% Options for packages loaded elsewhere
% Options for packages loaded elsewhere
\PassOptionsToPackage{unicode}{hyperref}
\PassOptionsToPackage{hyphens}{url}
\PassOptionsToPackage{dvipsnames,svgnames,x11names}{xcolor}
%
\documentclass[
  english,
  11pt,
]{article}
\usepackage{xcolor}
\usepackage[margin=25mm]{geometry}
\usepackage{amsmath,amssymb}
\setcounter{secnumdepth}{-\maxdimen} % remove section numbering
\usepackage{iftex}
\ifPDFTeX
  \usepackage[T1]{fontenc}
  \usepackage[utf8]{inputenc}
  \usepackage{textcomp} % provide euro and other symbols
\else % if luatex or xetex
  \usepackage{unicode-math} % this also loads fontspec
  \defaultfontfeatures{Scale=MatchLowercase}
  \defaultfontfeatures[\rmfamily]{Ligatures=TeX,Scale=1}
\fi
\usepackage{lmodern}
\ifPDFTeX\else
  % xetex/luatex font selection
\fi
% Use upquote if available, for straight quotes in verbatim environments
\IfFileExists{upquote.sty}{\usepackage{upquote}}{}
\IfFileExists{microtype.sty}{% use microtype if available
  \usepackage[]{microtype}
  \UseMicrotypeSet[protrusion]{basicmath} % disable protrusion for tt fonts
}{}
\makeatletter
\@ifundefined{KOMAClassName}{% if non-KOMA class
  \IfFileExists{parskip.sty}{%
    \usepackage{parskip}
  }{% else
    \setlength{\parindent}{0pt}
    \setlength{\parskip}{6pt plus 2pt minus 1pt}}
}{% if KOMA class
  \KOMAoptions{parskip=half}}
\makeatother
% Make \paragraph and \subparagraph free-standing
\makeatletter
\ifx\paragraph\undefined\else
  \let\oldparagraph\paragraph
  \renewcommand{\paragraph}{
    \@ifstar
      \xxxParagraphStar
      \xxxParagraphNoStar
  }
  \newcommand{\xxxParagraphStar}[1]{\oldparagraph*{#1}\mbox{}}
  \newcommand{\xxxParagraphNoStar}[1]{\oldparagraph{#1}\mbox{}}
\fi
\ifx\subparagraph\undefined\else
  \let\oldsubparagraph\subparagraph
  \renewcommand{\subparagraph}{
    \@ifstar
      \xxxSubParagraphStar
      \xxxSubParagraphNoStar
  }
  \newcommand{\xxxSubParagraphStar}[1]{\oldsubparagraph*{#1}\mbox{}}
  \newcommand{\xxxSubParagraphNoStar}[1]{\oldsubparagraph{#1}\mbox{}}
\fi
\makeatother


\usepackage{longtable,booktabs,array}
\newcounter{none} % for unnumbered tables
\usepackage{calc} % for calculating minipage widths
% Correct order of tables after \paragraph or \subparagraph
\usepackage{etoolbox}
\makeatletter
\patchcmd\longtable{\par}{\if@noskipsec\mbox{}\fi\par}{}{}
\makeatother
% Allow footnotes in longtable head/foot
\IfFileExists{footnotehyper.sty}{\usepackage{footnotehyper}}{\usepackage{footnote}}
\makesavenoteenv{longtable}
\usepackage{graphicx}
\makeatletter
\newsavebox\pandoc@box
\newcommand*\pandocbounded[1]{% scales image to fit in text height/width
  \sbox\pandoc@box{#1}%
  \Gscale@div\@tempa{\textheight}{\dimexpr\ht\pandoc@box+\dp\pandoc@box\relax}%
  \Gscale@div\@tempb{\linewidth}{\wd\pandoc@box}%
  \ifdim\@tempb\p@<\@tempa\p@\let\@tempa\@tempb\fi% select the smaller of both
  \ifdim\@tempa\p@<\p@\scalebox{\@tempa}{\usebox\pandoc@box}%
  \else\usebox{\pandoc@box}%
  \fi%
}
% Set default figure placement to htbp
\def\fps@figure{htbp}
\makeatother



\ifLuaTeX
\usepackage[bidi=basic,shorthands=off]{babel}
\else
\usepackage[bidi=default,shorthands=off]{babel}
\fi
\ifLuaTeX
  \usepackage{selnolig} % disable illegal ligatures
\fi


\setlength{\emergencystretch}{3em} % prevent overfull lines

\providecommand{\tightlist}{%
  \setlength{\itemsep}{0pt}\setlength{\parskip}{0pt}}



 


\makeatletter
\@ifpackageloaded{caption}{}{\usepackage{caption}}
\AtBeginDocument{%
\ifdefined\contentsname
  \renewcommand*\contentsname{Table of contents}
\else
  \newcommand\contentsname{Table of contents}
\fi
\ifdefined\listfigurename
  \renewcommand*\listfigurename{List of Figures}
\else
  \newcommand\listfigurename{List of Figures}
\fi
\ifdefined\listtablename
  \renewcommand*\listtablename{List of Tables}
\else
  \newcommand\listtablename{List of Tables}
\fi
\ifdefined\figurename
  \renewcommand*\figurename{Figure}
\else
  \newcommand\figurename{Figure}
\fi
\ifdefined\tablename
  \renewcommand*\tablename{Table}
\else
  \newcommand\tablename{Table}
\fi
}
\@ifpackageloaded{float}{}{\usepackage{float}}
\floatstyle{ruled}
\@ifundefined{c@chapter}{\newfloat{codelisting}{h}{lop}}{\newfloat{codelisting}{h}{lop}[chapter]}
\floatname{codelisting}{Listing}
\newcommand*\listoflistings{\listof{codelisting}{List of Listings}}
\makeatother
\makeatletter
\makeatother
\makeatletter
\@ifpackageloaded{caption}{}{\usepackage{caption}}
\@ifpackageloaded{subcaption}{}{\usepackage{subcaption}}
\makeatother
\usepackage{bookmark}
\IfFileExists{xurl.sty}{\usepackage{xurl}}{} % add URL line breaks if available
\urlstyle{same}
% fallback for those not using the hyperref driver hyperxmp:
\makeatletter
\@ifundefined{xmpquote}{\newcommand{\xmpquote}[1]{#1}}{}
\makeatother
\hypersetup{
  pdftitle={How a small wave grows{,} then fades},
  pdflang={en},
  colorlinks=true,
  linkcolor={blue},
  filecolor={Maroon},
  citecolor={Blue},
  urlcolor={Blue},
  pdfcreator={LaTeX via pandoc}}


\title{How a small wave grows, then fades}
\usepackage{etoolbox}
\makeatletter
\providecommand{\subtitle}[1]{% add subtitle to \maketitle
  \apptocmd{\@title}{\par {\large #1 \par}}{}{}
}
\makeatother
\subtitle{The background supplies energy; viscosity helps end the pulse}
\author{}
\date{2026-09-12}
\begin{document}
\maketitle


\href{../paper-map/article.html}{Start with the paper map} · Main route,
article 7 of 9

The previous article explained how waves can supply useful momentum
transport. But where do those waves get their energy? And how can they
disappear without leaving a large force when they are switched off?

The paper gives its pulses a carefully arranged life cycle. They begin
with a very small seed, grow by drawing energy from the background flow,
and later decay strongly enough that they can be cut off with small
errors.

The important competition is between \textbf{shear}, which can amplify a
disturbance, and \textbf{viscosity}, which smooths differences in
velocity.

\subsection{Begin with neighbouring lanes moving at different
speeds}\label{begin-with-neighbouring-lanes-moving-at-different-speeds}

Imagine parallel conveyor belts running at different speeds. A
disturbance that moves between them encounters different background
velocities.

A fluid with neighbouring layers moving at different speeds has shear.
That difference provides an energy source: an appropriately arranged
disturbance can draw energy from the background.

This does not mean every disturbance grows. Its direction and
relationship to the shear matter. Some motions extract energy; others do
not.

The paper designs the background and pulse orientations together to
obtain the needed amplification.

\subsection{See how the wave pattern
changes}\label{see-how-the-wave-pattern-changes}

Draw a stripe across several conveyor belts. As the belts move, they
carry different parts of the stripe different distances. The stripe
tilts.

A wave pattern in a sheared flow changes in an analogous way. Its
orientation evolves, and the spacing of its variations can become
smaller in the direction across the shear.

This matters because viscosity acts especially strongly on closely
spaced velocity differences.

\begin{figure}

\centering{

\pandocbounded{\includegraphics[keepaspectratio,alt={Initially aligned stripes tilt across a shear; a separate disturbance-speed curve rises to a peak and then decays.}]{figures/reader.png}}

}

\caption{\label{fig-pulse}The left panel illustrates stripes being
tilted by different background speeds. The right uses the companion's
simple straight-shear example: disturbance speed first grows and then
becomes small. This example illustrates transient amplification and
decay, not the exact rotating pulse mechanism of the paper.}

\end{figure}%

\subsection{Why fine patterns are smoothed
faster}\label{why-fine-patterns-are-smoothed-faster}

Think of temperature differences in a metal bar. Heat has only a short
distance to travel between neighbouring hot and cold bands when the
bands are tightly spaced.

Velocity differences also diffuse through viscosity. As a useful rule,
halving the variation length makes the diffusion time about four times
shorter. The distance enters squared.

You can therefore have an early stage when amplification wins and a
later stage when increasingly fine variations make damping win.

In the paper, the changing wave orientation also weakens amplification.
Both effects help produce growth followed by decay. The pulse cannot be
understood by comparing one fixed growth rate with one fixed damping
rate for its entire life.

\subsection{How the rotating background amplifies a
pulse}\label{how-the-rotating-background-amplifies-a-pulse}

Here is the source's physical picture in words. A fluctuation with
excess circular velocity can acquire an outward motion. At a larger
radius, the surrounding angular velocity is different. Under the chosen
shear conditions, the fluctuation's excess relative to its surroundings
can grow.

The effects reinforce one another. Axial shear provides an additional
route to amplification, especially where the rotational contribution
would be weak.

The paper makes these conditions quantitative and includes viscous
damping. Its selected rotating pulse model has an initial amplification
stage followed by decay.
\href{https://cdn.openai.com/pdf/32d9f210-8b73-45e0-91bc-82a30aef8a9a/navier-stokes.pdf}{Paper,
Section 2.2 and Lemma 7.4}

Our conveyor-belt picture explains shearing and energy exchange. It does
not reproduce every feature of this rotating instability.

\subsection{What the simpler computation
shows}\label{what-the-simpler-computation-shows}

The supporting experiment uses a known wave in a straight shear flow. We
can calculate its changing orientation and amplitude directly, then
compare with a numerical integration.

For its chosen dimensionless parameters, the disturbance speed rises to
about \textbf{3.6 times} its initial value before decaying. Since energy
depends on speed squared, the peak energy is about \textbf{13 times} its
initial value.

The gain comes from the background shear. It is not energy appearing
from nowhere, and it is not evidence that viscosity has become negative.

This straight-shear example demonstrates temporary amplification and
eventual viscous decay. The actual paper's pulse construction uses a
different, rotating and axially sheared setting. Keeping that
distinction lets the simpler example teach a mechanism without
pretending to reproduce the full construction.

\subsection{Why the quiet tails
matter}\label{why-the-quiet-tails-matter}

Imagine stopping a loudspeaker while its signal is large. An abrupt
cutoff introduces a sharp change. Turning it off after the signal has
already become tiny is easier to do gently.

The paper similarly applies temporal cutoffs in the pulses' small tails.
The errors introduced by the cutoff must be controlled, including their
space-time derivatives.

The seed force also has to remain compatible with the final smooth-force
requirement. The source arranges exponentially small seeds and tails;
merely drawing a bell-shaped amplitude curve would not establish those
estimates.

This connects the pulse's life cycle directly to the theorem's force
regularity. The pulse needs to be useful in the middle and sufficiently
quiet at its ends.

\subsection{Where this fits in the bigger
picture}\label{where-this-fits-in-the-bigger-picture}

\textbf{The wave's energy comes from the background, and its designed
decay makes its contribution controllable.} Averaged transport tells us
what the pulse does; amplification and damping tell us how it can do
that while respecting the viscous dynamics.

Adding pulses still creates new errors and can disturb important
weighted totals. Continue with
\href{../d6_moment_correction/article.html}{small local repairs that
preserve those totals}.

The \href{technical-article.html}{technical article},
\href{../../code/d5_shearing_wave/study.ipynb}{notebook}, and
\href{source-notes.md}{source notes} provide the exact straight-shear
example and its energy accounting.




\end{document}
